\documentclass[12pt,a4paper,oneside]{article}
\usepackage[utf8]{inputenc}
\usepackage[spanish,es-noshorthands]{babel}
\usepackage[T1]{fontenc}
\usepackage{lmodern}
\usepackage[margin=2.5cm]{geometry}
\usepackage{amsmath,amssymb}
\usepackage{xcolor}
\usepackage{hyperref}
\usepackage{graphicx}
\usepackage{booktabs}
\usepackage{tabularx}
\usepackage{enumitem}
\usepackage{multirow}
\usepackage{colortbl}
\usepackage{fancyhdr}
\usepackage{lastpage}
\usepackage{longtable}
\usepackage{listings}

\hypersetup{colorlinks=true,urlcolor=blue,linkcolor=black}

\definecolor{darkblue}{RGB}{31,56,100}
\definecolor{accent}{RGB}{79,143,247}
\definecolor{success}{RGB}{16,185,129}
\definecolor{warning}{RGB}{245,158,11}
\definecolor{danger}{RGB}{239,68,68}
\definecolor{lightgray}{RGB}{240,242,247}
\definecolor{codebg}{RGB}{248,249,252}
\definecolor{codekw}{RGB}{109,40,217}
\definecolor{codestr}{RGB}{16,120,40}
\definecolor{codecom}{RGB}{120,120,120}

\lstset{
  basicstyle=\ttfamily\small,
  backgroundcolor=\color{codebg},
  keywordstyle=\color{codekw}\bfseries,
  stringstyle=\color{codestr},
  commentstyle=\color{codecom}\itshape,
  showstringspaces=false,
  breaklines=true,
  frame=single,
  rulecolor=\color{lightgray},
  numbers=left,
  numberstyle=\tiny\color{codecom},
  xleftmargin=2em,
  framexleftmargin=1.5em
}

\pagestyle{fancy}
\fancyhf{}
\rhead{\textcolor{accent}{ISO SECURE}}
\lhead{\textcolor{accent}{Etapa 3 - Codificación e Integración}}
\cfoot{Página \thepage\ de \pageref{LastPage}}
\renewcommand{\headrulewidth}{0.5pt}
\renewcommand{\headrule}{\color{accent}\hrule}

\title{
  {\Huge \textbf{ISO SECURE}}\\
  {\LARGE Plataforma de Cumplimiento ISO 27001}\\[2ex]
  {\Large \textcolor{accent}{Etapa 3: Codificación e Integración}}\\[1ex]
  {\small Actividad Formativa 3 - 5\%}
}
\author{Agustín Peralta --- Universidad de San Buenaventura}
\date{Entrega: 11 de Mayo de 2026}

\begin{document}

\maketitle

\begin{abstract}
Este documento consolida las prácticas de codificación segura aplicadas en el proyecto ISO SECURE, el flujo de integración continua, la estrategia de control de versiones, los análisis estáticos y de dependencias, y las pruebas unitarias/integración que sustentan la fase de construcción del SDLC seguro. Se mapean las prácticas a OWASP ASVS 4.0 y a los controles del Anexo A de ISO/IEC 27001:2022 (cláusula A.8 — controles tecnológicos).
\end{abstract}

\tableofcontents
\newpage

% ============================================================================
\section{Introducción}

La Etapa 3 traduce las decisiones de arquitectura y diseño (Etapa 2) en código de producción siguiendo los principios del \textbf{Secure Software Development Lifecycle (SSDLC)} de OWASP SAMM y NIST SP 800-218. Esta fase es donde la mayor parte de las vulnerabilidades se introducen \textemdash{} y donde se previenen.

\textbf{Objetivos de la fase:}
\begin{itemize}[leftmargin=*]
  \item Codificar respetando las guías OWASP Secure Coding Practices.
  \item Establecer un flujo de Git con revisión obligatoria por pares.
  \item Automatizar análisis estático (SAST) y de dependencias (SCA) en CI.
  \item Garantizar cobertura mínima de pruebas (>70\% líneas críticas).
  \item Integrar los componentes (backend, frontend, base de datos) en entornos reproducibles.
\end{itemize}

% ============================================================================
\section{Prácticas de Codificación Segura}

\subsection{Mapeo a OWASP Top 10:2021}

\begin{longtable}{p{1cm}p{4cm}p{8cm}}
\toprule
\textbf{ID} & \textbf{Riesgo OWASP} & \textbf{Práctica aplicada en ISO SECURE} \\
\midrule
\endhead
A01 & Broken Access Control & \texttt{require\_role()} factory + helper \texttt{\_empresa\_filter} en todos los routers \\
A02 & Cryptographic Failures & TLS 1.3 obligatorio; secretos vía variables de entorno; tokens JWT firmados RS256 \\
A03 & Injection & ORM SQLAlchemy parametrizado; validación Pydantic v2 estricta; sin string concatenation en queries \\
A04 & Insecure Design & Modelado de amenazas previo (Etapa 2); revisión arquitectónica antes de codificar \\
A05 & Security Misconfiguration & Headers CSP, HSTS, X-Content-Type-Options en Nginx; \texttt{DEBUG=False} en producción \\
A06 & Vulnerable Components & Dependabot + \texttt{pip-audit} + \texttt{npm audit} en CI \\
A07 & Auth Failures & Supabase Auth gestiona contraseñas, MFA y bloqueo; tokens de corta vida \\
A08 & Data Integrity Failures & Hash de integridad en releases; verificación de firmas en webhooks n8n \\
A09 & Logging Failures & Logs estructurados JSON; mask de campos sensibles; envío a sink externo \\
A10 & SSRF & Whitelist de URLs salientes; sin proxy abierto; n8n con allowlist de hosts \\
\bottomrule
\end{longtable}

\subsection{Reglas internas de código seguro}

\textbf{Backend (Python):}
\begin{enumerate}[leftmargin=*]
  \item Nunca aceptar \texttt{empresa\_id} desde el body del request; inferirlo del usuario autenticado.
  \item Toda función pública en \texttt{routers/} declara dependencias \texttt{Depends()} explícitas (sesión + usuario).
  \item Schemas Pydantic separados por dirección: \texttt{XCreate}, \texttt{XUpdate}, \texttt{XResponse} (no reutilizar el modelo ORM como schema HTTP).
  \item Errores genéricos hacia el cliente en producción; detalles sólo en logs internos.
  \item \texttt{require\_role()} declarado al nivel del endpoint, no dentro de la función.
\end{enumerate}

\textbf{Frontend (React):}
\begin{enumerate}[leftmargin=*]
  \item Tokens en \texttt{localStorage} encapsulados en un adapter; no \texttt{innerHTML} ni \texttt{dangerouslySetInnerHTML}.
  \item Validación en cliente \textbf{además} de la del servidor; nunca como reemplazo.
  \item Interceptor Axios añade el JWT y captura 401 para redirigir al login.
  \item Variables de entorno con prefijo \texttt{VITE\_} y nunca incluir secretos en el bundle.
  \item Render condicional con guards de rol antes de renderizar acciones sensibles.
\end{enumerate}

\subsection{Ejemplos de implementación}

\textbf{Listado seguro de incidentes (backend):}

\begin{lstlisting}[language=Python]
@router.get("/incidents", response_model=List[IncidentResponse])
async def list_incidents(
    db: AsyncSession = Depends(get_db),
    user: UserProfile = Depends(get_current_user),
):
    # empresa_id se infiere del usuario, NUNCA del request
    empresa_id = _empresa_filter(user)
    return await incident_service.list_for(db, empresa_id=empresa_id)
\end{lstlisting}

\textbf{Guard de rol como factory (backend):}

\begin{lstlisting}[language=Python]
def require_role(allowed: list[str]):
    async def _guard(user: UserProfile = Depends(get_current_user)):
        if user.role not in allowed:
            raise HTTPException(403, detail="forbidden")
        return user
    return _guard

# Uso:
@router.put("/auth/role")
async def change_role(
    payload: ChangeRoleSchema,
    admin: UserProfile = Depends(require_role(["admin"])),
):
    ...
\end{lstlisting}

\textbf{Interceptor Axios (frontend):}

\begin{lstlisting}[language=JavaScript]
import axios from "axios";

const api = axios.create({
  baseURL: import.meta.env.VITE_API_URL + "/api/v1",
});

api.interceptors.request.use((config) => {
  const token = sessionStore.getToken();
  if (token) config.headers.Authorization = `Bearer ${token}`;
  return config;
});

api.interceptors.response.use(
  (r) => r,
  (err) => {
    if (err.response?.status === 401) sessionStore.logout();
    return Promise.reject(err);
  }
);
\end{lstlisting}

% ============================================================================
\section{Control de Versiones y Flujo Git}

\subsection{Estrategia de ramas}

Se adopta una variante de \textbf{Trunk-Based Development} adaptada al tamaño del equipo:

\begin{itemize}[leftmargin=*]
  \item \textbf{main}: rama protegida; sólo merges vía PR aprobado y CI verde.
  \item \textbf{feature/<slug>}: ramas de corta vida (<3 días); se mergean a main y se eliminan.
  \item \textbf{hotfix/<slug>}: parches urgentes desde main, mergeados con prioridad.
\end{itemize}

\subsection{Convención de commits}

Se usa \textbf{Conventional Commits} para habilitar changelog automático:

\begin{lstlisting}
<tipo>(<scope>): <descripcion>

tipos: feat, fix, docs, refactor, test, chore, perf, security
\end{lstlisting}

Ejemplos reales del repositorio:
\begin{itemize}
  \item \texttt{feat(etapa-1): Análisis e Ingeniería de Requerimientos - PDF entregable}
  \item \texttt{feat: agregar \_headers con CSP y cabeceras de seguridad}
  \item \texttt{docs: agregar metodologia PASTA + STRIDE para modelamiento de amenazas}
\end{itemize}

\subsection{Política de Pull Requests}

\begin{enumerate}[leftmargin=*]
  \item Toda PR debe pasar el pipeline de CI (lint, tests, SAST, SCA).
  \item Al menos un revisor humano aprueba antes del merge.
  \item No se mergean PRs con conflictos sin resolver.
  \item Squash & merge para mantener historial lineal.
  \item Los commits a main quedan bloqueados sin PR (branch protection).
\end{enumerate}

% ============================================================================
\section{Pipeline de Integración Continua}

\subsection{Etapas del pipeline}

\begin{center}
\begin{tabular}{|l|l|l|}
\hline
\rowcolor{accent!20}
\textbf{Etapa} & \textbf{Herramienta} & \textbf{Falla si} \\
\hline
Lint Python & ruff & errores E/F \\
\hline
Lint JS/JSX & ESLint + prettier & cualquier error \\
\hline
SAST Python & bandit & severidad >= medium \\
\hline
SAST JS & semgrep & regla bloqueante \\
\hline
SCA Python & pip-audit & CVE crítica/alta \\
\hline
SCA JS & npm audit + Snyk & CVE crítica \\
\hline
Tests unitarios & pytest + jest & cualquier fallo \\
\hline
Cobertura & coverage.py + jest --coverage & < 70\% backend \\
\hline
Build Docker & docker build & error de build \\
\hline
Scan imagen & Trivy & CVE crítica \\
\hline
\end{tabular}
\end{center}

\subsection{Configuración de ejemplo (GitHub Actions)}

\begin{lstlisting}[language=bash,morekeywords={name,on,jobs,runs,uses,steps,with}]
name: ci
on: [push, pull_request]
jobs:
  backend:
    runs-on: ubuntu-latest
    steps:
      - uses: actions/checkout@v4
      - uses: actions/setup-python@v5
        with: { python-version: "3.13" }
      - run: pip install -r requirements.txt -r requirements-dev.txt
      - run: ruff check .
      - run: bandit -r app -lll
      - run: pip-audit --strict
      - run: pytest --cov=app --cov-fail-under=70
  frontend:
    runs-on: ubuntu-latest
    steps:
      - uses: actions/checkout@v4
      - uses: actions/setup-node@v4
        with: { node-version: "20" }
      - working-directory: frontend
        run: |
          npm ci
          npm run lint
          npm audit --audit-level=high
          npm test -- --coverage
          npm run build
\end{lstlisting}

% ============================================================================
\section{Análisis Estático (SAST) y de Dependencias (SCA)}

\subsection{Herramientas seleccionadas}

\begin{table}[h]
\centering
\small
\begin{tabularx}{\textwidth}{lll}
\toprule
\textbf{Herramienta} & \textbf{Tipo} & \textbf{Cobertura} \\
\midrule
ruff & Linter & estilo + bugs comunes Python \\
bandit & SAST & vulnerabilidades comunes Python \\
semgrep & SAST multilenguaje & reglas custom + community rules \\
pip-audit & SCA Python & CVE en dependencias \\
npm audit & SCA Node & CVE registry \\
Trivy & Scanner contenedores & CVE en imagen Docker \\
Dependabot & Actualizaciones & PRs automáticas \\
\bottomrule
\end{tabularx}
\end{table}

\subsection{Hallazgos típicos y política de fix}

\begin{table}[h]
\centering
\small
\begin{tabularx}{\textwidth}{lXl}
\toprule
\textbf{Hallazgo} & \textbf{Acción} & \textbf{SLA} \\
\midrule
CVE crítica (CVSS >= 9.0) & Fix obligatorio antes del merge & 0 días \\
CVE alta (7.0-8.9) & Fix antes del próximo release & 7 días \\
CVE media (4.0-6.9) & Backlog priorizado & 30 días \\
Bandit MEDIUM/HIGH & Análisis manual + fix o supresión justificada & 14 días \\
ruff/eslint & Auto-fix vía pre-commit hook & 0 días \\
\bottomrule
\end{tabularx}
\end{table}

\subsection{Pre-commit hooks}

Se usa \texttt{pre-commit} con la siguiente configuración mínima para que los problemas se detecten antes del push:

\begin{lstlisting}[language=bash]
# .pre-commit-config.yaml
repos:
  - repo: https://github.com/astral-sh/ruff-pre-commit
    rev: v0.5.0
    hooks: [{ id: ruff, args: [--fix] }, { id: ruff-format }]
  - repo: https://github.com/PyCQA/bandit
    rev: 1.7.9
    hooks: [{ id: bandit, args: [-r, app] }]
  - repo: local
    hooks:
      - id: no-secrets
        name: detect-secrets
        entry: detect-secrets-hook
        language: system
\end{lstlisting}

% ============================================================================
\section{Estrategia de Pruebas}

\subsection{Pirámide de pruebas}

\begin{itemize}[leftmargin=*]
  \item \textbf{Base (70\%):} Unitarias — services y utils aislados.
  \item \textbf{Medio (20\%):} Integración — routers contra DB de prueba.
  \item \textbf{Top (10\%):} End-to-end — flujos críticos vía Playwright.
\end{itemize}

\subsection{Tipos de prueba}

\begin{table}[h]
\centering
\small
\begin{tabularx}{\textwidth}{lXl}
\toprule
\textbf{Tipo} & \textbf{Foco} & \textbf{Herramienta} \\
\midrule
Unitaria & Lógica pura de servicios, validaciones, helpers & pytest \\
Integración & Router + DB + Auth real & pytest-asyncio + httpx \\
Contrato & Compatibilidad de schemas Pydantic & schemathesis \\
Seguridad & RBAC, multi-tenancy, JWT inválido & pytest + scenarios \\
E2E & Login + listado + creación de incidente & Playwright \\
Carga & Throughput de /dashboard/summary & Locust o k6 \\
\bottomrule
\end{tabularx}
\end{table}

\subsection{Ejemplo: prueba de RBAC multi-tenant}

\begin{lstlisting}[language=Python]
@pytest.mark.asyncio
async def test_supervisor_solo_ve_su_empresa(client, db, empresas, users):
    # Arrange: 2 empresas con 1 incidente cada una
    sup_token = await login(users.supervisor_emp_a)

    # Act
    r = await client.get(
        "/api/v1/incidents",
        headers={"Authorization": f"Bearer {sup_token}"},
    )

    # Assert
    assert r.status_code == 200
    ids = {i["empresa_id"] for i in r.json()}
    assert ids == {str(empresas.a.id)}, "Filtro empresa_id falló"
\end{lstlisting}

\subsection{Cobertura mínima exigida}

\begin{itemize}[leftmargin=*]
  \item Backend global: $\geq$ 70\% líneas.
  \item Módulos críticos (auth, multi-tenant filter): $\geq$ 90\%.
  \item Frontend: $\geq$ 60\% líneas; 80\% en componentes con lógica de negocio.
\end{itemize}

% ============================================================================
\section{Integración de Componentes}

\subsection{Stack de integración}

Se usa \textbf{Docker Compose} para orquestar el entorno de desarrollo y staging:

\begin{lstlisting}[language=bash,morekeywords={services,build,ports,environment,depends_on,volumes}]
# docker-compose.yml (extracto real del repo)
version: "3.8"
services:
  api:
    build: .
    ports: ["8000:8000"]
    environment:
      - DATABASE_URL=${DATABASE_URL}
      - SUPABASE_URL=${SUPABASE_URL}
      - SUPABASE_ANON_KEY=${SUPABASE_ANON_KEY}
    depends_on: [postgres]
  postgres:
    image: postgres:15-alpine
    ports: ["5432:5432"]
    environment:
      - POSTGRES_PASSWORD=devpass
    volumes:
      - pgdata:/var/lib/postgresql/data
volumes:
  pgdata:
\end{lstlisting}

\subsection{Migración determinística}

El script \texttt{migrate.py} ejecuta los pasos de esquema en orden idempotente, evitando que dos despliegues lleguen a estados divergentes. Incluye:

\begin{enumerate}[leftmargin=*]
  \item Rename de enums legados (\texttt{partial $\rightarrow$ en\_proceso}).
  \item Creación de tipos enum nuevos.
  \item Creación de tablas (empresas, user\_profiles, cursos, etc.).
  \item Seed de cursos iniciales con URLs validadas.
  \item Adición de FK \texttt{empresa\_id} a incidents, controls, risk\_levels.
  \item Seed automático de los 93 controles ISO por sector detectado.
\end{enumerate}

\subsection{Entornos}

\begin{table}[h]
\centering
\begin{tabular}{lll}
\toprule
\textbf{Entorno} & \textbf{Propósito} & \textbf{Datos} \\
\midrule
dev local & Desarrollo personal & PostgreSQL en Docker + seed \\
staging & QA, demos, UAT & Supabase staging + datos sintéticos \\
producción & Usuarios reales & Supabase prod + backups diarios \\
\bottomrule
\end{tabular}
\end{table}

% ============================================================================
\section{Gestión de Secretos}

\begin{itemize}[leftmargin=*]
  \item \texttt{.env} en \texttt{.gitignore}; nunca commiteado.
  \item \texttt{.env.example} en el repo con placeholders.
  \item Producción: secretos en el gestor del proveedor (Supabase Vault, GitHub Actions Secrets).
  \item Rotación de la \texttt{SUPABASE\_ANON\_KEY} y \texttt{SERVICE\_ROLE\_KEY} cada 90 días.
  \item Hook \texttt{detect-secrets} en pre-commit bloquea pushes con tokens accidentales.
\end{itemize}

% ============================================================================
\section{Trazabilidad: Diseño $\rightarrow$ Implementación}

\begin{table}[h]
\centering
\small
\begin{tabularx}{\textwidth}{llX}
\toprule
\textbf{Decisión Etapa 2} & \textbf{Implementación} & \textbf{Validación} \\
\midrule
Filtro empresa\_id server-side & \texttt{\_empresa\_filter} en routers & Test \texttt{test\_multitenant\_isolation} \\
JWT validado vs Supabase & \texttt{auth.py::get\_current\_user} & Test \texttt{test\_jwt\_invalid\_rejected} \\
RBAC factory & \texttt{require\_role()} & Test \texttt{test\_supervisor\_no\_puede\_borrar} \\
ORM parametrizado & SQLAlchemy 2.0 async & SAST bandit sin findings \\
Migraciones idempotentes & \texttt{migrate.py} con guards \texttt{IF NOT EXISTS} & CI corre migrate dos veces \\
Logs sin secretos & Mask en \texttt{logging.Filter} & Grep en logs de staging \\
\bottomrule
\end{tabularx}
\end{table}

% ============================================================================
\section{Conclusiones}

La Etapa 3 materializa la arquitectura en código siguiendo prácticas de codificación segura mapeadas a OWASP Top 10:2021 y a OWASP ASVS 4.0. Los puntos clave del cierre de esta fase son:

\begin{enumerate}[leftmargin=*]
  \item Pipeline CI con SAST, SCA y tests es \textbf{bloqueante}: ningún merge sin CI verde.
  \item Cobertura mínima 70\% global y 90\% en módulos críticos (auth, multi-tenant).
  \item Convención Conventional Commits + branch protection en \texttt{main}.
  \item Pre-commit hooks llevan el shift-left al máximo: el desarrollador ve los problemas antes del push.
  \item Cada decisión arquitectónica de la Etapa 2 tiene su contrapartida verificable en código y prueba.
\end{enumerate}

La fase de codificación no termina con el primer release: las prácticas, hooks y políticas descritas son parte del ciclo de mejora continua exigido por ISO 27001:2022 (cláusulas 9 y 10).

\section*{Referencias}
\begin{itemize}[leftmargin=*]
  \item OWASP Top 10:2021.
  \item OWASP Application Security Verification Standard (ASVS) 4.0.3.
  \item OWASP Secure Coding Practices Quick Reference Guide v2.0.
  \item NIST SP 800-218: Secure Software Development Framework (SSDF).
  \item ISO/IEC 27001:2022 — Anexo A.8 (Controles tecnológicos).
  \item Conventional Commits 1.0.0 — \url{https://www.conventionalcommits.org/}.
\end{itemize}

\end{document}
